% ===========================================================================
\section{What a self-naming rate can and cannot show}
\label{sec:measure}

\paragraph{Terms.} A \emph{panel} is a set of models that both write the
texts and judge them, so that each judge has texts of its own among those it
judges. For judge $J$:
\begin{itemize}
\item the \textbf{self-naming rate} (hit rate) is the share of $J$'s own
texts that $J$ attributes to $J$;
\item the \textbf{false-alarm rate} is the share of texts $J$ did not write
that $J$ attributes to $J$;
\item the \textbf{peer baseline} is the share of $J$'s texts that the
\emph{other} judges attribute to $J$, i.e.\ how identifiable $J$'s writing is
to anyone else under the same protocol;
\item \textbf{self-advantage} is the self-naming rate minus the peer
baseline.
\end{itemize}
The first two are the hit and false-alarm rates of signal detection theory
\citep{green1966signal, macmillan2005detection}: they separate a judge's
ability to tell its own text from others' (sensitivity) from its readiness to
answer ``me'' (bias). The peer baseline adds what signal detection alone
does not: a reference for how attributable the text is. Self-advantage is
positive only if a judge does better on its own text than judges who did not
write it. It is not sufficient on its own: a judge that names itself on half
of everything it reads beats peers who rarely name it. A judge therefore
counts as self-recognising only if it also \emph{discriminates}, with a hit
rate above its false-alarm rate (tested by permuting which text in each
lineup is its own).

\paragraph{Three real cases.} All three come from our frontier panel
(Section~\ref{sec:frontier}). In the lineup, Claude names itself on 33 of its
38 texts (87\%), on only 14 of the 152 texts it did not write (9\%), and the
other four judges name Claude on 82 of their 152 judgments of Claude's texts
(54\%). It discriminates, and it beats its peers by 33 points: this is
self-recognition. Shown the same texts one at a time, Claude again names
itself on 33 of 38 with 13 false alarms, but now its peers name Claude on 125
of 152 (82\%), because in that format Claude receives over half of all
guesses, whoever wrote the text. The self-naming rate is unchanged; the evidence for
self-recognition is almost gone (+5 points). In the same format GPT names
itself on 36 of its 38 texts (95\%), which looks like the strongest
recognition in the study, but it also names itself on 79 of the 152 texts it
did not write (52\%): much of that rate is a readiness to say ``me''. The
self-naming rate alone cannot tell these three cases apart.

\paragraph{Which statistics we report.} For every judge and format we give
the self-naming rate with an exact 95\% interval, the false-alarm rate, the
peer baseline and the self-advantage with a prompt-level bootstrap interval.
Judgments on the same prompt are not independent (five judges see the same
five texts), so every test resamples or permutes whole prompts or whole
responses, never single judgments; Appendix~\ref{app:stats} gives the
effective sample sizes. For yes/no questions, which yield a probability of
``yes'' rather than a name, we report the area under the ROC curve (AUC) of
that probability for own versus others' texts, which is unaffected by a
judge's overall readiness to say yes.
